%!tex

\magnification=1400

\input ../../../TeX/macros/macros.tex
\input tikz

\usetikzlibrary{arrows.meta, math}

\fig{%
  \tikzpicture[scale=3.5]
    \def\err(#1,#2){
      \def\errt{.01}
      \def\errl{.05}
      \draw (#1-\errt,#2) -- (#1+\errt,#2);
      \draw [Bracket-Bracket] (#1,#2-\errl) -- (#1,#2+\errl);
    }
    \err(0.587, 0.0596988)
    \err(1.076, 0.1235132)
    \err(1.565, 0.1790230)
    \err(2.054, 0.2482254)
    \err(2.543, 0.3180960)
    \draw  (0,-0.019570847407407294) -- (3.13,0.3910458705354843);
    \draw [dashed, draw=black!30!white] (0,0) -- (0,0.3910458705354843) node [above] {$\lambda^2$}; 
    \draw [dashed, draw=black!30!white] (0,0) -- (3.13,0) node [right] {$T$};
  \endtikzpicture
}{%
  Regress\~ao linear de $\lambda^2$ em $T$ (tra\c c\~ao).
}

\advance\proclcounter by 4
\procl{Quest\~ao}{}{}
Via script {\tt ondas-estacionarias.py} determina-se que o coeficiente $B$ do termo dependente \'e
$$
  B=0.13118744982201008.
$$

Agora, visto que
$$
  \lambda^2={1\over \mu f^2}\cdot T,
$$
logo,
$$
  B={1\over \mu f^2},
$$
donde
$$
  \mu={1\over Bf^2}\approx0.0005293528042405269
$$

aproximando para o primeiro d\'\i gito significativo, obtemos
$$
  \mu\approx0.0005\,{\rm kg}/{\rm m}
$$

\procl{Quest\~ao}{}{}
Via c\'alculo indireto temos
$$
  \mu_{\rm ind}={m\over\ell}={0.98\,{\rm kg}\over1.592\,{\rm m}}\approx0.0006155778894472361.
$$
Por outro lado, via experimental
$$
  \mu_{\rm exp}\approx0.0005293528042405269.
$$
Pela teoria de propraga\c c\~ao de erros tem-se
$$
  \eqalign{
    \delta_\mu &=\sqrt{({\partial_B\mu}\cdot\sigma_y)^2}\cr
               &=\sigma_B|\partial_B\mu|\cr
               &=\sigma_B\Bigl({1\over Bf}\Bigr)^2\cr
               &\approx0.00001.\cr
  }
$$
\'E consp\'\i cuo que 
$$
  \mu_{\rm ind}=0.0006\notin[0.00052,0.00054]
$$

N\~ao, os resultados n\~ao s\~ao compat\'\i veis.






\bye
